\subsection{Diagonalizable algebras and etale algebras}

DEFINITION 1. --- Let A be an algebra over K; then A is said to be diagonalizable if there exists an integer \(n \geqslant 0\) such that A is isomorphic to the product algebra \(\mathbf{K}^n\) . We say that A is diagonalized by an extension L of K if the algebra \(\mathbf{A}_{(L)}\) over L derived from A by extension of scalars is diagonalizable. We shall say that A is etale if there exists an extension of K which diagonalizes A.

We recall that the product algebra \(K^{n}\) is the vector space \(K^{n}\) equipped with the product defined by

\[
\left(x _ {1}, \dots , x _ {n}\right) \cdot \left(y _ {1}, \dots , y _ {n}\right) = \left(x _ {1} y _ {1}, \dots , x _ {n} y _ {n}\right).\tag{5}
\]

If \(\varepsilon_1, \ldots, \varepsilon_n\) is the canonical basis of \(\mathbf{K}^n\) , we have

\[
\varepsilon_ {i} ^ {2} = \varepsilon_ {i}, \quad \varepsilon_ {i} \varepsilon_ {j} = 0 \quad \text { if } \quad i \neq j\tag{6}
\]

and \(1 = \varepsilon_{1} + \dots +\varepsilon_{n}\)

Every etale algebra over K is commutative and of finite degree over K.

PROPOSITION 1. --- Let A be an algebra of finite degree n over the field K ; then the following conditions are equivalent :

\begin{enumerate}
\def\labelenumi{\alph{enumi})}
\item
  The algebra \(\mathbf{A}\) is diagonalizable.
\item
  There is a basis \((e_1, \ldots, e_n)\) of A such that \(e_i^2 = e_i\) and \(e_i e_j = 0\) for \(i \neq j\) .
\item
  The K-algebra homomorphisms of A into K generate the dual of the vector K-space A.
\item
  Every A-module is a direct sum of submodules which are of dimension 1 over K.
\end{enumerate}

The equivalence of \(a)\) and \(b)\) follows from Formula (6); on the other hand the \(n\) projections \(\mathbf{K}^n \to \mathbf{K}\) are algebra homomorphisms, hence \(a)\) implies \(c)\) . If \(c)\) holds, the algebra homomorphisms of \(\mathbf{A}\) into \(\mathbf{K}\) form a basis of the dual of \(\mathbf{A}\) (V, p.~27, Th. 1), we denote them by \(u_1, \ldots, u_n\) ; then \(a \mapsto (u_i(a))\) is an isomorphism of \(\mathbf{A}\) onto the algebra \(\mathbf{K}^n\) , whence \(a)\) . We have thus established the equivalence of the conditions \(a), b)\) and \(c)\) .

Suppose that \(b)\) holds and let \(\mathbf{M}\) be an A-module; then the homotheties \((e_i)_{\mathbf{M}}\) of ratio \(e_i\) are projectors of \(\mathbf{M}\) , and \(\mathbf{M}\) is a direct sum of the \(e_i\mathbf{M}\) , which are sub-A-modules. We may thus suppose that there exists an index \(i\) such that \((e_j)_{\mathbf{M}} = 0\) for \(j \neq i\) . Hence every vector subspace of \(\mathbf{M}\) is a sub-A-module, whence \(d)\) .

Conversely suppose that \(d\) holds and consider the A-module \(\mathbf{A}_{s}\) . There exists then a basis \((f_{i})\) of the vector K-space \(\mathbf{A}\) such that \(\mathbf{A}f_{i} = \mathbf{K}f_{i}\) for \(i = 1, \ldots, n\) . After replacing each \(f_{i}\) by a suitable scalar multiple, if necessary, we may suppose that \(1 = f_{i} + \cdots + f_{n}\) . If \(i \neq j\) , then \(f_{i}f_{j}\) belongs to \(\mathbf{A}f_{i} \cap \mathbf{A}f_{j} = \mathbf{K}f_{i} \cap \mathbf{K}f_{j}\) hence it is zero. Now \(f_{i} = f_{i}f_{1} + \cdots + f_{i}f_{n} = f_{i}^{2}\) , whence \(b\) ).

COROLLARY. --- Let L be an extension of K and H the set of algebra homomorphisms of A into L. We have Card \(H \leq [A:K]\) , with equality if and only if A is diagonalized by L. If A is diagonalized by L, then H is a basis of the vector L-space \(\operatorname{Hom}_{\mathbf{K}}(\mathbf{A},\mathbf{L})\) .

The vector space \(\operatorname{Hom}_{\mathbf{K}}(\mathbf{A},\mathbf{L})\) over L has dimension \([A:K]\) , by Formula (2), and H is a free subset of \(\operatorname{Hom}_{\mathbf{K}}(\mathbf{A},\mathbf{L})\) by Th. 1 (V, p.~27). We thus have Card \(H \leqslant [A:K]\) with equality if and only if H is a basis of \(\operatorname{Hom}_{\mathbf{K}}(\mathbf{A},\mathbf{L})\) . There exists an isomorphism of vector L-spaces, say \(\pi: \operatorname{Hom}_{\mathbf{K}}(\mathbf{A}, \mathbf{L}) \to \mathbf{A}_{(\mathbf{L})}^{*}\) , characterized by \(u(x) = (\pi u)(1 \otimes x)\) for \(x \in \mathbf{A}\) , and \(\pi\) maps \(\mathcal{H}\) onto the set \(\mathcal{H}_{\mathbf{L}}\) of L-algebra homomorphisms of \(\mathbf{A}_{(\mathbf{L})}\) into L. Finally the equivalence of \(a)\) and \(c)\) in Prop. 1 shows that the algebra \(\mathbf{A}_{(\mathbf{L})}\) over L is diagonalizable if and only if \(\mathcal{H}_{\mathbf{L}}\) generates the vector space \(\mathbf{A}_{(\mathbf{L})}^{*}\) over L. This completes the proof of the Corollary.

PROPOSITION 2. --- Let A be an algebra over K and \(\Omega\) an algebraically closed extension of K. The following assertions are equivalent:

\begin{enumerate}
\def\labelenumi{\alph{enumi})}
\item
  The algebra \(\mathbf{A}\) is etale.
\item
  There exists an extension of finite degree which diagonalizes A.
\item
  The extension \(\Omega\) of \(\mathbf{K}\) diagonalizes \(\mathbf{A}\) .
\end{enumerate}

Suppose that A is etale. Let n be the degree of A over K, let L be an extension of K which diagonalizes A and let H be the set of algebra homomorphisms of A into L. By the Cor. to Prop. 1 we have Card H = n.~On the other hand, for each \(u \in H\) , we have \([u(A):K] \leqslant n\) . By V, p.~18, Th. 2, the subextension \(L'\) of L generated by the images of elements of H is of finite degree over K. Since there exist n distinct homomorphisms of A into \(L'\) , the extension \(L'\) diagonalizes A, by the Cor. 1 of Prop. 1. This shows that a) implies b).

Since every extension of finite degree of K is isomorphic to a subextension of \(\Omega\) (V, p.~20, Th. 1), b) implies c). Finally c) clearly implies a).

